\documentclass[11pt,a4paper]{article}
\usepackage[utf8]{inputenc}
\usepackage[vietnamese]{babel}
\usepackage{amsmath,amssymb}
\usepackage{geometry}
\geometry{a4paper, total={180mm,260mm}, left=15mm, right=15mm, top=18mm, bottom=18mm}

\usepackage{tikz}
\usetikzlibrary{shapes.geometric, arrows.meta, positioning, calc}
\tikzset{
  arr/.style={-{Stealth[scale=0.8]}},
  every node/.append style={align=center}
}
\usepackage{tcolorbox}
\tcbuselibrary{skins,breakable}
\usepackage{fontawesome}
\usepackage{fancyhdr}
\usepackage{enumitem}

\pagestyle{fancy}
\fancyhf{}
\renewcommand{\headrulewidth}{0.5pt}
\renewcommand{\footrulewidth}{0.5pt}
\fancyhead[L]{\small\bfseries\color{violet!85!black}\faBook\ SỔ TAY KIẾN THỨC TOÁN 2}
\fancyhead[R]{\small\bfseries\color{violet!85!black}Lớp Toán Thầy Lĩnh -- loptoanthaylinh.vn}
\fancyfoot[C]{\small\color{gray}-- Sổ tay kiến thức: Bài 32. Luyện tập chung thời gian --}

\setlength{\parindent}{0pt}
\setlength{\parskip}{0pt}

\begin{document}

% --- BANNER BÀI HỌC ---
\begin{tcolorbox}[
  enhanced,
  colback=violet!8!white,
  colframe=violet!70!black,
  arc=3mm, boxrule=1.2pt,
  left=3mm, right=3mm, top=2mm, bottom=2mm,
  halign=center
]
  {\LARGE\bfseries\color{violet!90!black}\faBookmark\ BÀI 32: LUYỆN TẬP CHUNG THỜI GIAN}\\[1mm]
  {\itshape\footnotesize\color{black!75} Đại tổng kết kiến thức Ngày, Giờ, Phút, Ngày, Tháng và lịch tuần lễ trên 1 trang A4}
\end{tcolorbox}

\vspace{-2mm}

\noindent
\begin{minipage}[t]{0.49\textwidth}\vspace{0pt}
\begin{tcolorbox}[
  enhanced,
  colback=violet!4!white,
  colframe=violet!75!black,
  arc=2mm, boxrule=0.8pt,
  left=2mm, right=2mm, top=1.8mm, bottom=1.8mm,
  title={\footnotesize\bfseries\color{white}\faKey\ 1. Bảng quy đổi thời gian}
]
{\footnotesize
\begin{itemize}\setlength{\itemsep}{1pt}\setlength{\leftskip}{-2.5mm}
  \item $1\text{ ngày} = 24\text{ giờ}$.
  \item $1\text{ giờ} = 60\text{ phút}$.
  \item $1\text{ tuần lễ} = 7\text{ ngày}$.
  \item $1\text{ năm} = 12\text{ tháng}$.
  \item Các tháng có $31$ ngày: $1, 3, 5, 7, 8, 10, 12$.
  \item Các tháng có $30$ ngày: $4, 6, 9, 11$.
  \item Tháng $2$: Có $28$ hoặc $29$ ngày.
\end{itemize}
}
\end{tcolorbox}
\end{minipage}
\hfill
\begin{minipage}[t]{0.49\textwidth}\vspace{0pt}
\begin{tcolorbox}[
  enhanced,
  colback=teal!4!white,
  colframe=teal!75!black,
  arc=2mm, boxrule=0.8pt,
  left=2mm, right=2mm, top=1.8mm, bottom=1.8mm,
  title={\footnotesize\bfseries\color{white}\faSliders\ 2. Kĩ năng giải toán thời gian}
]
{\footnotesize
\begin{itemize}\setlength{\itemsep}{1pt}\setlength{\leftskip}{-2.5mm}
  \item \textbf{Tính ngày hôm qua, hôm nay, ngày mai:}
  \begin{itemize}\setlength{\itemsep}{0.5pt}\setlength{\leftskip}{-2.5mm}
    \item Hôm qua: Bớt $1$ ngày.
    \item Ngày mai: Thêm $1$ ngày.
  \end{itemize}
  \item \textbf{Cộng trừ tuần lễ:}\\
  Thêm $7$ ngày hoặc bớt $7$ ngày thì thứ trong tuần không thay đổi.
  \item \textbf{So sánh khoảng thời gian:}\\
  $1\text{ giờ} > 50\text{ phút}$ (vì $1\text{ giờ} = 60\text{ phút}$).
\end{itemize}
}
\end{tcolorbox}
\end{minipage}

\vspace{-1.5mm}
\begin{tcolorbox}[
  enhanced,
  colback=orange!4!white,
  colframe=orange!80!black,
  arc=2mm, boxrule=0.8pt,
  left=2mm, right=2mm, top=1.5mm, bottom=1.5mm,
  title={\footnotesize\bfseries\color{white}\faSitemap\ 3. Sơ đồ tư duy: Trục tuần hoàn các đơn vị thời gian}
]
\centering
\begin{tikzpicture}[x=1.4cm, y=0.8cm]
  \node[rectangle, rounded corners=2pt, draw=violet!80!black, fill=violet!15, font=\scriptsize\bfseries] (phut) at (0,0) {PHÚT\\($60\text{ phút} = 1\text{ h}$)};
  \node[rectangle, rounded corners=2pt, draw=teal!80!black, fill=teal!15, font=\scriptsize\bfseries] (gio) at (2.5,0) {GIỜ\\($24\text{ giờ} = 1\text{ ngày}$)};
  \node[rectangle, rounded corners=2pt, draw=orange!80!black, fill=orange!15, font=\scriptsize\bfseries] (ngay) at (5,0) {NGÀY\\($7\text{ ngày} = 1\text{ tuần}$)};
  \node[rectangle, rounded corners=2pt, draw=green!80!black, fill=green!15, font=\scriptsize\bfseries] (thang) at (7.5,0) {THÁNG\\($12\text{ tháng} = 1\text{ năm}$)};

  \draw[->, thick, violet!80!black] (phut) -- (gio);
  \draw[->, thick, teal!80!black] (gio) -- (ngay);
  \draw[->, thick, orange!80!black] (ngay) -- (thang);
\end{tikzpicture}
\end{tcolorbox}

\vspace{-1.5mm}
\noindent
\begin{minipage}[t]{0.49\textwidth}\vspace{0pt}
\begin{tcolorbox}[
  enhanced,
  colback=green!4!white,
  colframe=green!60!black,
  arc=2mm, boxrule=0.8pt,
  left=2mm, right=2mm, top=1.8mm, bottom=1.8mm,
  title={\footnotesize\bfseries\color{white}\faCheckCircle\ 4. Ví dụ thực hành chuẩn mực}
]
{\footnotesize
\textbf{1. Điền số thích hợp:}\\
- $1$ ngày $2$ giờ $= 24 + 2 = \mathbf{26\text{ giờ}}$.\\
- $1$ tuần lễ và $3$ ngày $= 7 + 3 = \mathbf{10\text{ ngày}}$.
\textbf{2. Bài toán lịch:}\\
Nếu thứ Năm tuần này là ngày $10$, thì thứ Năm tuần trước là ngày mấy?\\
\textit{Giải:} Thứ Năm tuần trước là ngày: $10 - 7 = \mathbf{3}$.
}
\end{tcolorbox}
\end{minipage}
\hfill
\begin{minipage}[t]{0.49\textwidth}\vspace{0pt}
\begin{tcolorbox}[
  enhanced,
  colback=red!4!white,
  colframe=red!75!black,
  arc=2mm, boxrule=0.8pt,
  left=2mm, right=2mm, top=1.8mm, bottom=1.8mm,
  title={\footnotesize\bfseries\color{white}\faExclamationTriangle\ 5. Mẹo tránh bẫy trắc nghiệm}
]
{\footnotesize
\begin{itemize}\setlength{\itemsep}{1pt}\setlength{\leftskip}{-2.5mm}
  \item \textbf{Bẫy qua tháng mới:}\\
  Nếu hôm nay là ngày $31$ tháng $5$, thì ngày mai là ngày \textbf{$1$ tháng $6$}, không có ngày $32$!
  \item Tuần sau $\implies (+7\text{ ngày})$. Tuần trước $\implies (-7\text{ ngày})$.
\end{itemize}
}
\end{tcolorbox}
\end{minipage}


\end{document}
